%% generated by gen_tables.py - RE-ANALYSIS 2026
\begin{table}[htbp]
\centering
\small
\caption{Air properties used by the analytical model (at the inlet and outlet bulk temperatures)}
\label{tab:air}
\begin{tabular}{lrrl}
\toprule
Property & 300 K & 368.9 K & Unit \\
\midrule
Density $\rho = p_{op}/(RT)$ & 1.1766 & 0.9569 & kg/m$^3$ \\
Specific heat $c_p$ & 1007 & 1011 & J/(kg\,K) \\
Dynamic viscosity $\mu$ & $1.846\times10^{-5}$ & $2.176\times10^{-5}$ & Pa\,s \\
Thermal conductivity $k$ & 0.02624 & 0.03123 & W/(m\,K) \\
Prandtl number $Pr$ & 0.707 & 0.696 & -- \\
\bottomrule
\end{tabular}
\par\smallskip\parbox{\linewidth}{\footnotesize\raggedright Source: \texttt{02\_Engineering\_Calculations/MATERIAL\_PROPERTIES.md}; $c_p$, $\mu$, $k$, $Pr$ from Incropera et al.\ Table A.4 \cite{incropera}; density computed from the ideal-gas law at 101,325~Pa with $R = 287.058$~J/(kg\,K) because tabulated densities refer to 100~kPa. In Fluent, $c_p$, $\mu$ and $k$ are piecewise-linear over 8 points between 250 and 600~K.}
\end{table}
